% !TEX root = ../main.tex
%% Sujet type bac n° 2 — sujet inédit OptiTAB (valeurs vérifiées : verifications/sujets_bac.py)

\sujetBac{2}{{Le dihydrogène vert : électrolyse, gaz parfait, pile à combustible}/9,{L'échographie Doppler : ultrasons, atténuation, effet Doppler}/6,{Ötzi : datation au carbone 14}/5}

\begin{exoBac}{9}{Le dihydrogène vert}
Produit par électrolyse de l'eau à l'aide d'électricité d'origine renouvelable, le dihydrogène dit « vert » permet de
stocker de l'énergie, puis de la restituer sous forme électrique dans une pile à combustible. On étudie les trois
étapes de cette chaîne : production, stockage et restitution.
\begin{donnees}
Constante de Faraday : $F=\qty{96500}{C.mol^{-1}}$ ; constante des gaz parfaits : $R=\qty{8,314}{J.mol^{-1}.K^{-1}}$ ;
$M(\ce{H2})=\qty{2,0}{g.mol^{-1}}$ ; $T(\unit{K})=\theta(\unit{\degreeCelsius})+273$ ; \qty{1}{bar} $=$ \qty{1e5}{Pa} ; \qty{1}{kW.h} $=$ \qty{3,6e6}{J}.
Couples oxydant-réducteur : $\ce{H+}/\ce{H2}$ et $\ce{O2}/\ce{H2O}$.
\end{donnees}

\partieExo{1}{Produire le dihydrogène}
\begin{docu}{1}{L'électrolyseur à membrane}
Dans chaque cellule, l'eau est décomposée en dihydrogène et en dioxygène ; une membrane laisse passer les ions
$\ce{H+}$ d'un compartiment à l'autre. Une cellule fonctionne sous la tension $U=\qty{1,9}{V}$, traversée par un courant
d'intensité $I=\qty{200}{A}$.
\schema{bac-s2-electrolyseur}
\end{docu}
\begin{questions}
\item Écrire les demi-équations électroniques des réactions aux électrodes, puis l'équation de l'électrolyse
$\ce{2H2O -> 2H2 + O2}$. Pourquoi cette transformation est-elle qualifiée de forcée ?
\item Identifier, en justifiant, l'électrode où se forme le dihydrogène et indiquer le sens de circulation des électrons
dans les fils de connexion.
\item Calculer la quantité d'électrons échangés en une heure de fonctionnement d'une cellule, puis la masse de
dihydrogène produite.
\item Calculer le volume occupé par ce dihydrogène à \qty{20}{\degreeCelsius} sous \qty{1,0}{bar}, en le considérant comme un gaz parfait.
\item Calculer l'énergie électrique consommée par la cellule en une heure. Montrer que la production d'un kilogramme
de dihydrogène demande environ \qty{51}{kW.h}.
\end{questions}

\partieExo{2}{Stocker le dihydrogène}
\begin{docu}{2}{La bouteille}
Le dihydrogène est comprimé dans une bouteille de \qty{50}{L}, sous \qty{200}{bar}, à \qty{20}{\degreeCelsius}. À cette pression, le comportement
du dihydrogène s'écarte du modèle du gaz parfait : la masse réellement stockée est 1,13 fois plus faible que celle
calculée avec ce modèle.
\end{docu}
\begin{questions}[resume]
\item Calculer, avec le modèle du gaz parfait, la masse de dihydrogène contenue dans la bouteille.
\item Pourquoi le modèle du gaz parfait est-il mis en défaut à haute pression ? Calculer la masse réellement stockée.
\end{questions}

\partieExo{3}{Restituer l'énergie}
\begin{docu}{3}{La pile à combustible}
La pile consomme le dihydrogène de la bouteille, soit \qty{0,73}{kg}, et le dioxygène de l'air ; elle ne rejette que de
l'eau. Sa tension de fonctionnement vaut \qty{0,70}{V}.
\end{docu}
\begin{questions}[resume]
\item Écrire les demi-équations des réactions à l'anode et à la cathode de la pile, puis l'équation de sa réaction de
fonctionnement.
\item Calculer la charge électrique débitée par la pile, puis l'énergie électrique qu'elle fournit, en kilowattheures.
\item \initiative Évaluer le rendement de la chaîne « électrolyse, stockage, pile à combustible ». Commenter ce résultat.
\end{questions}
\end{exoBac}

\begin{exoBac}{6}{L'échographie Doppler}
Pour mesurer la vitesse du sang dans une artère du cou, un médecin applique une sonde sur la peau. La sonde émet des
ultrasons de fréquence $f=\qty{5,0}{MHz}$, réfléchis par les globules rouges, puis reçoit l'écho.
\schema{bac-s2-doppler}
\begin{docu}{1}{Le décalage Doppler}
Lorsque les globules rouges se déplacent à la vitesse $v$, le décalage entre la fréquence reçue et la fréquence émise
vaut $\Delta f=f_\mathrm{R}-f=\dfrac{2\,f\,v\cos\theta}{c}$ ; il est positif lorsque le sang se rapproche de la sonde.
\end{docu}
\begin{docu}{2}{L'atténuation des ultrasons}
Les tissus absorbent les ultrasons : l'atténuation vaut \qty{0,5}{dB} par centimètre parcouru et par mégahertz de
fréquence. L'appareil ne détecte plus un écho atténué de plus de \qty{60}{dB}.
\end{docu}
\begin{donnees}
Célérité des ultrasons dans les tissus : $c=\qty{1540}{m.s^{-1}}$ ; $\theta=\ang{60}$ ; profondeur de l'artère : $d=\qty{4,0}{cm}$.
Atténuation en décibels : $10\log\dfrac{I_\text{émise}}{I_\text{reçue}}$.
\end{donnees}
\begin{questions}
\item Les ultrasons sont-ils des ondes mécaniques ou électromagnétiques ? Pourquoi l'oreille ne les perçoit-elle pas ?
\item Calculer la longueur d'onde des ultrasons dans les tissus.
\item Calculer la durée qui sépare l'émission d'une salve et la réception de l'écho renvoyé par l'artère.
\item L'appareil mesure $\Delta f=+\qty{1,3}{kHz}$. Le sang se rapproche-t-il de la sonde ? Calculer sa vitesse.
\item Calculer l'atténuation de l'écho dû à l'absorption par les tissus, puis le rapport des intensités émise et reçue.
\item Pourquoi incline-t-on la sonde plutôt que de la placer perpendiculairement à l'artère ?
\item \initiative Pour examiner le cœur, situé à une dizaine de centimètres sous la peau, le médecin utilise une sonde de
\qty{2,5}{MHz}. Justifier ce choix, puis expliquer pourquoi on utilise une fréquence plus élevée lorsque c'est possible.
\end{questions}
\end{exoBac}

\begin{exoBac}{5}{Ötzi, l'homme des glaces}
En 1991, des randonneurs découvrent dans un glacier des Alpes le corps momifié d'un homme, aussitôt surnommé Ötzi. Les
scientifiques datent sa mort à l'aide du carbone 14.
\begin{docu}{1}{La datation au carbone 14}
Tant qu'un organisme est vivant, il échange du carbone avec son environnement : l'activité du carbone 14 y reste égale
à $A_0=\qty{0,23}{Bq}$ par gramme de carbone. Après la mort, cette activité décroît : $A(t)=A_0\,\mathrm{e}^{-\lambda t}$. La demi-vie du
carbone 14 vaut $t_{1/2}=\num{5700}$ ans. Les tissus d'Ötzi ont une activité de \qty{0,122}{Bq} par gramme de carbone.
\end{docu}
\begin{donnees}
Noyaux : $\ce{^{14}_{6}C}$, $\ce{^{14}_{7}N}$ ; une année vaut \qty{3,16e7}{s} ; $\lambda=\dfrac{\ln2}{t_{1/2}}$ ; l'activité est liée au nombre $N$ de noyaux
radioactifs par $A=\lambda\,N$.
\end{donnees}
\begin{questions}
\item Le carbone 14 est radioactif $\beta^-$. Écrire l'équation de sa désintégration en précisant les lois de conservation
utilisées.
\item Définir la demi-vie. Calculer la constante radioactive $\lambda$ en $\text{an}^{-1}$, puis en $\unit{s^{-1}}$.
\item Montrer que la durée écoulée depuis la mort vaut $t=\dfrac{1}{\lambda}\ln\dfrac{A_0}{A}$.
\item \initiative Depuis combien de temps Ötzi est-il mort ?
\item Calculer le nombre de noyaux de carbone 14 contenus dans un gramme de carbone d'un organisme vivant.
\item Pourquoi cette méthode ne permet-elle pas de dater des restes de plus de \num{50000} ans ?
\end{questions}
\end{exoBac}
